% ---
% id: meat-core/req:gradient-at-maximal-support
% kind: req
% title: "Gradient evaluation at maximal support"
% summary: >-
%   The engine evaluates a gradient of the pinched trace power only at an
%   iterate whose support is kept maximal, for instance after an explicit
%   perturbation. At lower support a direction that enlarges the support can
%   have an infinite one-sided derivative ([[thm:gradient]] (ii)).
% prominence: standard
% status: implemented
% uses: [thm:gradient, lem:branch-support]
% mentions: [def:branch-support, meat-core/frank-wolfe/lem:fo-face, meat-core/frank-wolfe/rem:fo-nopert, meat-core/frank-wolfe/req:first-order-withdrawal]
% assumes: []
% refs: []
% provenance: newly-derived
% implements: certiqs_key.meat_core.objective.gradient:deficit_gradient
% lean:
% ---

Theorem~\ref{thm:gradient}\,(ii) places every infinite one-sided
derivative at a direction that strictly enlarges the support, and at
maximal support no feasible direction does (Lemma~\ref{lem:fo-face}). A
gradient taken after an explicit perturbation is the gradient of another
function (Remark~\ref{rem:fo-nopert}), and on the first-order path it
removes that iterate's bracket from the bound
(Requirement~\ref{req:first-order-withdrawal}).

\begin{requirement}\label{req:gradient-at-maximal-support}
The engine evaluates a gradient of the pinched trace power only at an
iterate whose support is (kept) maximal, maximality being that of
\eqref{eq:Wcs} for the branch operators, for instance after an explicit
perturbation.
\end{requirement}
